\documentclass[a4paper,12pt]{article}
\usepackage[utf8]{inputenc}
\usepackage[magyar]{babel}
\usepackage{booktabs}
\usepackage{geometry}
\usepackage{xcolor}
\usepackage{titlesec}
\usepackage{graphicx}
\usepackage{subcaption}
\usepackage{hyperref}

\geometry{margin=25mm}

\titleformat{\section}{\large\bfseries\color{darkblue}}{\thesection}{1em}{}
\definecolor{darkblue}{rgb}{0.0, 0.2, 0.4}

\title{\textbf{Tor Website Fingerprinting\\Önálló laboratórium beszámoló}}
\author{Pálos Ferenc}
\date{\today}

\begin{document}

\maketitle

\section*{Összefoglaló}
A beszámoló a Tor hálózaton végzett weboldal ujjlenyomat-készítés (Website Fingerprinting, WF)
feladatát mutatja be, két eltérő modellcsalád összehasonlításával: klasszikus gépi tanulás
(Random Forest) és mélytanulás (Triplet MLP + KNN). A vizsgálat a zárt világú
(baseline, obfs4), a zero-shot és a nyílt világú (open-world) forgatókönyvekben méri a
teljesítményt. A végső értékelés a futtatott Python pipeline által számolt metrikákon és a
generált ábrákon alapul. A Random Forest erős baseline és obfs4 teljesítményt ad, míg a
Triplet MLP a zero-shot beállításokban kedvezőbb általánosítást mutat. A dokumentum
részletesen bemutatja az adatgyűjtési és reprodukálhatósági lépéseket is, hogy a
pipeline végponttól végpontig újrafuttatható legyen.

\section{Célkitűzések és hozzájárulások}
A projekt célja egy reprodukálható WF pipeline kialakítása és két modellcsalád objektív
összehasonlítása a Tor forgalomban jelentkező domain shift és ismeretlen forgalom kezelésére.
Ennek keretében:
\begin{itemize}
    \item a PCAP állományoktól a kiértékelési metrikákig egy végig automatizált feldolgozási láncot raktam össze,
    \item aggregált folyamjellemzőket definiáltam és egységes 21-dimenziós feature teret alkalmaztam,
    \item \textit{shared} és \textit{optimized} top-10 feature track-eket vezettem be,
    \item egységes ábragenerálást biztosítottam a kiértékelési JSON-okból,
    \item az open-world ``Other'' osztály kezelését és a zero-shot beállításokat külön mérési szálként kezeltem.
\end{itemize}

\section{Adatok és forgatókönyvek}
A mérés három adatdoménen történik: baseline Tor forgalom, obfs4 Tor forgalom, valamint
az open-world ``Other'' kategória. A fő forgatókönyvek:
\begin{itemize}
    \item \textbf{Baseline (closed-world):} standard Tor böngészés.
    \item \textbf{Obfs4 (closed-world):} obfs4 pluggable transport által elfedett forgalom.
    \item \textbf{Zero-shot:} baseline-on tanított modellek tesztelése obfs4 forgalmon.
    \item \textbf{Open-world:} ``Other'' kategória hozzáadása az ismeretlen oldalak modellezésére.
\end{itemize}
A zárt világú osztályokat rétegzett 80\%/20\% tanító/teszt felosztással értékelem. Az
open-world ``Other'' osztályt ezzel szemben 50\%--50\% arányban bontom tanító/teszt
részre, mert itt a cél nem az adott ``Other'' webhelyek pontos felismerése, hanem az
ismeretlen forgalomra való általánosítás. A nagyobb ``Other'' teszt rész valósághűbb
``ismeretlen'' mintaállományt ad, csökkenti a recall becslés zaját, és akkor is
értelmezhető méretű tesztet hagy, ha az ``Other'' mintaszám kisebb.

\section{Adatgyűjtés és előfeldolgozás}
Az adatgyűjtés automatizált headless böngészéssel és \texttt{tcpdump}-pal történik, aktív Tor
SOCKS proxy és ControlPort mellett. A SOCKS proxy a Tor hálózatra irányítja a böngésző
forgalmát, míg a ControlPort a Tor példány vezérlését teszi lehetővé (pl. új áramkör
kérése, állapotlekérdezés). A zárt világú baseline/obfs4 forgalmat a
\texttt{scripts/collection/collector.py} (futtató: \texttt{run\_collector.sh}) gyűjti, míg
az open-world ``Other'' adatokat a \texttt{server\_collect\_other.py} szkript (futtató:
\texttt{run\_server\_collect\_other.sh}) állítja elő és a \texttt{tor\_dataset/other\_tor/}
könyvtárba menti. A gyűjtés környezetét opcionálisan környezeti változók (pl.
\texttt{TOR\_WF\_PCAP\_DIR}, \texttt{TOR\_WF\_INTERFACE}, \texttt{TOR\_WF\_GUARD\_IP},
\texttt{TOR\_WF\_CAPTURE\_DURATION}, \texttt{TOR\_WF\_WARMUP\_DURATION}) szabályozzák.
Itt a \texttt{TOR\_WF\_INTERFACE} adja meg a \texttt{tcpdump} által figyelt hálózati
interfészt (pl. \texttt{eth0} vagy \texttt{wlan0}), a \texttt{TOR\_WF\_WARMUP\_DURATION}
pedig a tényleges mérés előtti bemelegítési időt másodpercben, amely a Tor/circuit
stabilizálását szolgálja.

\section{Feldolgozási pipeline és jellemzők}
A pipeline három lépésből áll. \textbf{(1)} A PCAP fájlokból kiolvasom az egyes csomagokat,
és csomagonkénti idősorokat készítek (mikor jött a csomag, mekkora volt, milyen irányban
haladt). \textbf{(2)} Ezekből az idősorokból egy-egy kapcsolatot összefoglaló, aggregált
jellemzőket számolok. Ilyenek például az összes be- és kimenő csomagszám, a be- és kimenő
bájtok összege, a csomagok közti érkezési idők statisztikái (átlag, szórás, kvantilisek),
a kapcsolat időtartama, illetve az irányváltások száma. \textbf{(3)} A jellemzőkön tanítom
és értékelem a modelleket.
Az alap 21 jellemzőből minden futásban top-10 készletet választok, így a modellek összehasonlítása
egységes dimenziószámon történik.
A jellemzők főbb csoportjait a \ref{tab:feature_groups} táblázat foglalja össze.

\begin{table}[hbt!]
    \caption{Aggregált folyamjellemzők csoportjai és példák.}
    \vspace{2mm}
    \centering
    \begin{tabular}{ll}
        \toprule
        \textbf{Csoport} & \textbf{Jellemzők (példák)} \\ \midrule
        Csomagszámok & total\_pkts, in\_pkts, out\_pkts \\
        Bájtstatisztikák & total\_bytes, in\_bytes, out\_bytes, in/out ratio \\
        Időköz-statisztikák & iat\_mean, iat\_std, iat\_median, iat\_p90 \\
        Méretstatisztikák & size\_mean, size\_std, size\_median, size\_p90 \\
        Időtartam és dinamika & duration, direction\_changes \\
        \bottomrule
    \end{tabular}
    \label{tab:feature_groups}
\end{table}

\section{Modellek és kiértékelés}
\textbf{Random Forest.} 300 döntési fából álló modell, stabilizált top-10 jellemzőkiválasztással
és három véletlen maggal (42, 52, 62). A három mag segítségével a tanító/teszt felosztást
és a modell tanítását háromszor ismétlem, majd átlagot és szórást számolok.\\
\textbf{Triplet MLP + KNN.} A Triplet MLP 32 dimenziós beágyazást tanul. A Triplet Margin Loss
arra kényszeríti a modellt, hogy az azonos weboldalhoz tartozó minták beágyazása közel legyen,
míg a különböző weboldalaké távol. A beágyazásokon $k=5$ KNN osztályozást végzek.\\
\textbf{Feature track-ek.} \textit{shared} track-ben egy közös top-10 készletet választok ki a
mutual information alapján. \textit{optimized} track-ben forgatókönyvenként külön top-10
készlet készül (RF: több futásból stabilizált fontossági rangsor, DL: mutual information).\\
\textbf{Mélytanulási beállítások.} A Triplet MLP 30 epochon tanul standardizált bemeneten,
BatchNorm és Dropout rétegekkel, hogy csökkentse a túltanulást.

\section{Kiértékelési metrikák}
A teljesítményt pontossággal és macro-F1 értékkel mérem. A pontosság a helyes predikciók
aránya, a macro-F1 pedig az osztályonként számolt F1-értékek átlaga:
\begin{equation}
    \mathrm{Acc} = \frac{1}{N} \sum_{i=1}^{N} \mathbf{1}[\hat{y}_i = y_i], \quad
    \mathrm{MacroF1} = \frac{1}{C} \sum_{c=1}^{C} \mathrm{F1}_c.
\end{equation}
Open-world esetben külön vizsgálom az ``Other'' osztály recall értékét is, mivel ez mutatja,
hogy az ismeretlen forgalmat mennyire képes a modell leválasztani.

\section{Ábrák és vizuális összegzés}
A \texttt{generate\_all\_figures.py} szkript a mentett metrikákból generálja a végső ábrákat,
ezért az ábrák és a táblázatok ugyanabból a kiértékelési kimenetből készülnek.
Az ábrák célja, hogy a fő trendeket gyorsan összegezze, majd részletesebb bontást adjon
forgatókönyv, feature track és osztályszintű viselkedés szerint.

\begin{figure}[hbt!]
    \centering
    \includegraphics[width=0.92\linewidth]{../figures/01_closed_world_summary.png}
    \caption{Zárt világú teljesítmény-összegzés (pontosság és macro-F1) a két modellre.}
    \label{fig:closed_world_summary}
\end{figure}

\begin{figure}[hbt!]
    \centering
    \includegraphics[width=0.92\linewidth]{../figures/02_macro_f1_all_scenarios.png}
    \caption{Macro-F1 alakulása minden forgatókönyvben, zárt és nyílt világú esetekben.}
    \label{fig:macro_f1_all}
\end{figure}

\begin{figure}[hbt!]
    \centering
    \includegraphics[width=0.92\linewidth]{../figures/10_accuracy_all_scenarios.png}
    \caption{Pontosság alakulása minden forgatókönyvben, zárt és nyílt világú esetekben.}
    \label{fig:accuracy_all}
\end{figure}

\begin{figure}[hbt!]
    \centering
    \includegraphics[width=0.92\linewidth]{../figures/03_shared_vs_optimized_delta.png}
    \caption{A \textit{shared} és \textit{optimized} track-ek közti különbségek (delta) pontosságban és macro-F1-ben.}
    \label{fig:shared_vs_optimized_delta}
\end{figure}

\begin{figure}[hbt!]
    \centering
    \includegraphics[width=0.92\linewidth]{../figures/04_open_world_focus.png}
    \caption{Nyílt világú forgatókönyvek összehasonlítása.}
    \label{fig:open_world_focus}
\end{figure}

\begin{figure}[hbt!]
    \centering
    \includegraphics[width=0.92\linewidth]{../figures/05_feature_importance_shared.png}
    \caption{A top-10 jellemzők fontossága a közös (shared) kiválasztás mellett.}
    \label{fig:feature_importance_shared}
\end{figure}

\begin{figure}[hbt!]
    \centering
    \includegraphics[width=0.92\linewidth]{../figures/06_feature_importance_separate.png}
    \caption{Forgatókönyvenként optimalizált top-10 jellemzők összehasonlítása.}
    \label{fig:feature_importance_separate}
\end{figure}

\begin{figure}[hbt!]
    \centering
    \includegraphics[width=0.92\linewidth]{../figures/11_feature_stability_rf.png}
    \caption{A top-10 jellemzők stabilitása Random Forest esetén (átlag $\pm$ szórás a több futásból).}
    \label{fig:feature_stability_rf}
\end{figure}

\begin{figure}[hbt!]
    \centering
    \begin{subfigure}[b]{0.49\linewidth}
        \centering
        \includegraphics[width=\linewidth]{../figures/07_confusion_matrices.png}
        \caption{Alap konfúziós mátrixok.}
        \label{fig:confusion_matrices_raw}
    \end{subfigure}
    \hfill
    \begin{subfigure}[b]{0.49\linewidth}
        \centering
        \includegraphics[width=\linewidth]{../figures/08_confusion_matrices_recall.png}
        \caption{Recall-normalizált változatok.}
        \label{fig:confusion_matrices_recall}
    \end{subfigure}
    \caption{Konfúziós mátrixok a főbb forgatókönyvekben (shared track, seed=42).}
    \label{fig:confusion_matrices_pair}
\end{figure}

\begin{figure}[hbt!]
    \centering
    \includegraphics[width=0.85\linewidth]{../figures/09_other_recall_open_world.png}
    \caption{Az ``Other'' osztály recall értékei az open-world forgatókönyvekben.}
    \label{fig:other_recall_open_world}
\end{figure}

\section{Kísérleti beállítások}
A tanítás és kiértékelés rétegzett tanító/teszt felosztással történik. A zárt világú
osztályoknál 80\%/20\% arányt használok, míg az open-world ``Other'' osztályt 50\%--50\%
arányban bontom az ismeretlen forgalomra való általánosítás hangsúlyozására. A stabilitást
három véletlen mag biztosítja (42, 52, 62). Az eredményeket átlag $\pm$ szórás formában
közlöm. A metrikák a pontosság (accuracy) és a macro-F1. A zero-shot beállításban a baseline
adatokon tanított modelleket obfs4 teszt adatokon értékelem. A táblázatok a \textit{shared}
track eredményeit mutatják.

\section{Reprodukálhatóság és futtatás}
A pipeline fő lépései az alábbi parancsokkal futtathatók:
\begin{itemize}
    \item \texttt{python src/extract\_all\_features.py}
    \item \texttt{python src/evaluate\_all\_rf.py}
    \item \texttt{python src/evaluate\_all\_dl.py}
    \item \texttt{python src/generate\_all\_figures.py}
\end{itemize}

\section{Eredmények}
\begin{table}[hbt!]
    \caption{Zárt világú (Standard) forgatókönyvek eredményei.}
    \vspace{2mm}
    \centering
    \begin{tabular}{lcccc}
        \toprule
        \textbf{Forgatókönyv} & \textbf{RF Acc.} & \textbf{RF Macro-F1} & \textbf{DL Acc.} & \textbf{DL Macro-F1} \\ \midrule
        Baseline & 80.44$\pm$0.96 & 79.76$\pm$0.99 & 61.19$\pm$2.22 & 60.23$\pm$2.47 \\
        Obfs4 & 73.33$\pm$1.78 & 71.88$\pm$1.57 & 50.64$\pm$2.59 & 49.48$\pm$2.70 \\
        Zero-Shot & 15.08$\pm$1.37 & 10.83$\pm$1.57 & 23.01$\pm$2.59 & 18.54$\pm$2.64 \\
        \bottomrule
    \end{tabular}
\end{table}

\begin{table}[hbt!]
    \caption{Nyílt világú (Open-World) forgatókönyvek eredményei.}
    \vspace{2mm}
    \centering
    \begin{tabular}{lcccc}
        \toprule
        \textbf{Forgatókönyv} & \textbf{RF Acc.} & \textbf{RF Macro-F1} & \textbf{DL Acc.} & \textbf{DL Macro-F1} \\ \midrule
        OW-Baseline & 80.93$\pm$0.84 & 79.11$\pm$0.68 & 62.00$\pm$1.32 & 60.37$\pm$0.76 \\
        OW-Obfs4 & 74.10$\pm$1.59 & 71.03$\pm$1.65 & 48.90$\pm$2.58 & 48.46$\pm$1.78 \\
        OW-Zero-Shot & 22.73$\pm$1.79 & 10.41$\pm$1.00 & 27.27$\pm$1.79 & 18.42$\pm$0.41 \\
        \bottomrule
    \end{tabular}
\end{table}

\section{Elemzés és megvitatás}
A Random Forest zárt világú baseline és obfs4 környezetben egyértelműen magasabb
pontosságot és macro-F1-t ér el, ami a kézzel tervezett aggregált jellemzők hatékony
kihasználását jelzi (lásd \ref{fig:closed_world_summary} és \ref{fig:accuracy_all}).
A zero-shot beállításban mindkét modell teljesítménye drasztikusan romlik, különösen
open-world esetben, ami a domain shift egyik legfontosabb gyakorlati kihívását mutatja
(lásd \ref{fig:macro_f1_all} és \ref{fig:open_world_focus}). A Triplet MLP ugyanakkor
stabilabb zero-shot értékeket ad (Standard Zero-Shot: DL 23.01\% vs. RF 15.08\%), ami a
beágyazások általánosíthatóságára utal.
Az \textit{optimized} track-ek több forgatókönyvben javítást hoznak, de a javulás nem
egyenletes, ezért a \textit{shared} beállítás jó kompromisszumot ad a konszisztencia és
teljesítmény között (lásd \ref{fig:shared_vs_optimized_delta}). A stabil top-10 lista
azt sugallja, hogy a leginformatívabb aggregált jellemzők robusztusak a tanító/teszt
felosztás változására (lásd \ref{fig:feature_stability_rf}).
Az open-world forgatókönyvekben a recall-normalizált konfúziós mátrixok jól kiemelik,
hogy az ``Other'' osztályt a modellek gyakran sikeresen külön választják a zárt világú
osztályoktól, ugyanakkor ez nem feltétlenül jelent magas precíziót
(lásd \ref{fig:confusion_matrices_recall} és \ref{fig:other_recall_open_world}).

\section{Korlátok}
\begin{itemize}
    \item A vizsgálat aggregált folyamjellemzőkre épül, ami részben elfedheti a finom időbeli mintázatokat.
    \item Az open-world ``Other'' osztály sokféle weboldal forgalmát fogja össze, ezért a modell
    nehezebben választja el a zárt világú osztályoktól, ami csökkentheti a precíziós mutatókat.
    \item A top-10 jellemzők statikus kiválasztása miatt a ritkább, de informatív mintázatok elveszhetnek.
\end{itemize}

\section{Következtetések és jövőbeli munka}
A projekt igazolja, hogy a klasszikus és mélytanulási modellek eltérő erősségekkel rendelkeznek
a Tor WF feladatban. A Random Forest magas baseline teljesítményt ad, míg a Triplet MLP
zero-shot helyzetekben jobb általánosítást mutat. A jövőben a drift-robosztus tanítás
(az időben változó forgalmi eloszláshoz igazított tréning, pl. időablakos validálás vagy
domain-invariáns reprezentációk) és az adaptív újratanulás (drift detektálás után
periodikus vagy online frissítés friss adatokból) lehet indokolt. Emellett a finomabb
idősoros jellemzők integrálása is ígéretes (csomag-idősor, burst-mintázatok, n-gram vagy
embedding alapú leírás), mert ezek érzékenyebben ragadják meg a forgalom dinamikáját.
Végül érdemes a modellfrissítés költség-hatékonyságát is vizsgálni: mennyi új adat és
számítási idő szükséges adott teljesítmény-nyereséghez, és milyen frissítési gyakoriság
adja a legjobb kompromisszumot.

\begin{thebibliography}{9}
    \bibitem{cherubin2022}
    Cherubin, G., et al. (2022). \textit{A study on concept drift in website fingerprinting}.
    \bibitem{sirinam2019}
    Sirinam, P., et al. (2019). \textit{Deep Fingerprinting: Undermining Website Fingerprinting Defenses with Deep Learning}.
\end{thebibliography}

\end{document}